This paper presented a replication of the work by Hou et al.\ \cite{hou2021rampmeter},
who applied reinforcement learning to freeway ramp metering signal control
and compared it against three conventional controllers.
The replication aimed to verify the claim that the reinforcement learning algorithms
Ape-X DQN, PPO and A3C outperform the baseline controllers
no ramp metering, fixed-time control and ALINEA
under mixed near-congested/congested and extremely congested traffic conditions.

\medskip

After implementing the environment and controllers as described in the original paper,
the obtained results did not match those reported by Hou et al.,
particularly for the two metrics we consider most important, 
traffic throughput and average vehicle speed.
With a deviation of up to 13.03\% in average vehicle speed for the no ramp meter case (mixed traffic scenario),
we tuned the simulation parameters empirically to bring the baseline values closer to the reported ones.

\medskip

With the tuned parameters, the deviation in average vehicle speed was reduced to at most 2.53\% and the deviation in traffic throughput to at most 10.19\%
across the baseline controllers. The remaining metrics, however, still showed notable deviations
and the claimed performance gains of the reinforcement learning controllers could not be observed, as
all three RL algorithms achieved similar performance to the baselines in both the mixed and the extreme traffic scenario.
The reward curves of the RL algorithms also did not converge to the values reported in the original paper.
One observation that we can confirm is that the training of A3C was significantly slower
than that of Ape-X DQN and PPO.

\medskip

Because the original paper does not specify the exact location and technical configuration of the measurements,
several assumptions had to be made when implementing the evaluation metrics.

\medskip

Overall, the replication revealed that the results reported by Hou et al.\ could not be reproduced
based on the information available in the original paper.
The claimed performance gains of the reinforcement learning controllers
over the baseline controllers could not be confirmed.
A more complete specification of the simulation environment and the measurement procedure would be necessary to enable a fair and conclusive comparison
between reinforcement-learning-based and classical ramp metering strategies.
